% Contenido del póster: Serfes, Wu y Avramidis (2025), Journal of Banking & Finance. DOI 10.1016/j.jbankfin.2024.107338
% Elaborado con apoyo de IA a partir del resumen oficial del artículo (OpenAlex), 2026-11-12.
\def\DOI{10.1016/j.jbankfin.2024.107338}
\def\TituloEs{FinTech frente a banco: el impacto de la tecnología crediticia en la competencia del mercado de crédito}
\def\TituloOriginal{FinTech vs. bank: The impact of lending technology on credit market competition}
\def\Autores{Serfes, Wu y Avramidis}
\def\Anio{2025}
\def\Revista{Journal of Banking \& Finance}

\def\Problema{La tecnología se ha extendido en la evaluación de créditos. La pregunta es si eso cambia la competencia en el mercado de préstamos a empresas y qué tipo de crédito ofrece cada prestamista.}
\def\Objetivo{Explicar cómo compiten las fintech y los bancos cuando difieren en su capacidad de evaluar a los prestatarios y en las garantías que exigen.}
\def\Metodo{Modelo teórico con prestamistas que tienen distinta capacidad de evaluación (\textit{screening}) y prestatarios con proyectos de distinto horizonte de inversión. Dos prestamistas inicialmente simétricos compiten usando la tecnología de evaluación y las exigencias de garantía.}
\def\Resultados{\begin{itemize}
  \item Las fintech ofrecen sobre todo préstamos sin garantía a prestatarios con proyectos de corto plazo.
  \item Los bancos tradicionales ofrecen sobre todo préstamos respaldados por activos a proyectos de largo plazo.
  \item Los prestamistas usan la evaluación y las garantías para reducir la competencia, lo que segmenta el mercado y restringe la oferta de crédito.
  \item Cuando mejora la tecnología de información, el efecto sobre la oferta de crédito y las tasas depende del segmento.
\end{itemize}}
\def\Conclusion{La segmentación entre crédito fintech sin garantía de corto plazo y crédito bancario con garantía de largo plazo ofrece una explicación, desde la oferta, del crecimiento del crédito sin garantía.}
\def\Relevancia{Este resultado justifica comparar el financiamiento participativo peruano con las tasas tradicionales de hasta 360 días. Como la SMV no publica el plazo de cada operación, ese tramo es un supuesto declarado. En el artículo del autor, la TEA bancaria de pequeña y microempresa de hasta 360 días promedió \cifra{teaBcosPyme2022}~\% en 2022 y \cifra{teaBcosPyme2026}~\% en 2026.}
\def\PalabrasClave{fintech, competencia bancaria, garantías, evaluación crediticia, segmentación del mercado}
% --- Ficha del estudio (herramientas/posteres/ficha_articulos.py) ---
\def\Ambito{Teórico}
\def\Enfoque{Modelo de competencia con evaluación crediticia y garantías}
\def\FiguraPeru{3}
% --- fin de la ficha ---
